\documentclass[../DoAn.tex]{subfiles}
\begin{document}

Bốn chương trước đã trình bày đầy đủ quá trình từ phân tích yêu cầu, lựa chọn
công nghệ đến thiết kế, hiện thực và kiểm thử hệ thống ALIBARBIE. Chương này
tổng hợp ba đóng góp kỹ thuật đặc thù mà đồ án đánh giá là có giá trị thực tiễn
cao nhất, không chỉ trong phạm vi dự án này mà còn có thể áp dụng cho các hệ
thống thương mại điện tử tương tự. Mỗi đóng góp được trình bày theo cấu trúc
ba phần: (i) dẫn dắt bài toán, (ii) giải pháp thiết kế và hiện thực, và (iii)
kết quả đạt được.

%%%%%%%%%%%%%%%%%%%%%%%%%%%%%%%%%%%%%%%%%%
\section{Quy trình hoàn trả hàng với vòng đời trạng thái rõ ràng}
\label{section:5.1}

\subsection{Bài toán}
\label{subsection:5.1.1}

Hoàn trả hàng là nghiệp vụ đặc thù của thương mại điện tử, nhưng phần lớn
các nền tảng SaaS và mã nguồn mở hiện tại xử lý tính năng này theo hướng đơn
giản hóa: khách hàng gửi yêu cầu, nhân viên xem xét và ra quyết định chấp
nhận hoặc từ chối. Mô hình hai trạng thái này bỏ qua thực tế vận hành phức
tạp hơn nhiều: khi nhận yêu cầu, nhân viên cần liên hệ khách hàng để xác minh
thông tin; tiếp theo cần theo dõi việc khách hàng gửi hàng về; sau khi nhận
được hàng vật lý, quản lý mới có cơ sở để kiểm tra và ra quyết định; cuối
cùng, việc hoàn tiền cần được tách biệt khỏi quyết định chấp nhận để đảm bảo
tính kiểm soát tài chính.

Nếu không mô hình hóa rõ ràng các bước này, hệ thống dễ dẫn đến các vấn đề
nghiêm trọng trong vận hành: nhân viên khác nhau không biết yêu cầu đang ở
giai đoạn nào; không có ai chịu trách nhiệm rõ ràng ở từng bước; không thể
thống kê được bao nhiêu yêu cầu đang chờ nhận hàng hay đang chờ xử lý hoàn
tiền. Ngoài ra, nếu nghiệp vụ chỉ do quản lý xử
lý toàn bộ mà không tách phân quyền theo từng bước, hệ thống không đảm bảo
kiểm soát nội bộ.

Bài toán đặt ra là thiết kế một mô hình trạng thái cho quy trình hoàn trả hàng
đáp ứng đồng thời ba yêu cầu: phản ánh đúng vòng đời thực tế của một yêu cầu
hoàn trả từ lúc phát sinh đến khi kết thúc; phân định rõ vai trò và quyền hạn
của từng tác nhân tại từng bước; và đảm bảo không có chuyển đổi trạng thái bất
hợp lệ có thể xảy ra do lỗi thao tác hoặc lỗi kỹ thuật.

\subsection{Giải pháp}

\subsubsection*{Mô hình máy trạng thái hữu hạn}
Giải pháp thiết kế quy trình hoàn trả hàng dưới dạng \textbf{máy trạng thái
hữu hạn} (Finite State Machine --- FSM) với bảy trạng thái và các chuyển đổi
hợp lệ được kiểm soát chặt chẽ. Bảng \ref{tab:return-states} mô tả ý nghĩa
và tác nhân chịu trách nhiệm tại từng trạng thái.

\begin{table}[H]
\centering
\caption{Bảng trạng thái vòng đời yêu cầu hoàn trả hàng}
\label{tab:return-states}
\begin{tabular}{|p{4cm}|p{4cm}|p{5cm}|}
\hline
\textbf{Trạng thái} & \textbf{Tác nhân chịu trách nhiệm} & \textbf{Điều kiện chuyển sang trạng thái tiếp theo} \\
\hline
\texttt{PENDING\_REVIEW} & Hệ thống tạo tự động & Nhân viên xem xét và bắt đầu liên hệ \\
\hline
\texttt{CONTACTING} & Nhân viên & Xác nhận thông tin xong, hướng dẫn khách gửi hàng \\
\hline
\texttt{WAITING\_ITEM} & Nhân viên & Nhân viên nhận được hàng trả về \\
\hline
\texttt{ITEM\_RECEIVED} & Nhân viên & Quản trị viên kiểm tra hàng và ra quyết định \\
\hline
\texttt{APPROVED} & Quản trị viên & Quản trị viên xử lý hoàn tiền \\
\hline
\texttt{REJECTED} & Quản trị viên & Trạng thái kết thúc \\
\hline
\texttt{REFUNDED} & Quản trị viên & Trạng thái kết thúc \\
\hline
\end{tabular}
\end{table}

Các chuyển đổi hợp lệ được biểu diễn theo đồ thị có hướng như sau:

\[
\begin{aligned}
\texttt{PENDING\_REVIEW} 
&\to \texttt{CONTACTING} \\
&\to \texttt{WAITING\_ITEM} \\
&\to \texttt{ITEM\_RECEIVED} \to
\begin{cases}
\texttt{APPROVED} \to \texttt{REFUNDED} \\
\texttt{REJECTED}
\end{cases}
\end{aligned}
\]

Chỉ có các chuyển đổi được liệt kê trên là hợp lệ. Mọi yêu cầu chuyển đổi
không thuộc tập hợp này đều bị từ chối ở tầng nghiệp vụ.

\subsubsection*{Hiện thực kiểm soát chuyển đổi trong ControllerReturn}
Phương thức \texttt{UpdateStatus} trong \texttt{ControllerReturn} hiện thực
hóa logic kiểm soát chuyển đổi trạng thái thông qua hai cơ chế:

\textbf{Cơ chế 1: Bảng chuyển đổi hợp lệ.} Server định nghĩa tĩnh một bảng
ánh xạ từ trạng thái hiện tại sang tập các trạng thái có thể chuyển sang.
Khi nhận yêu cầu cập nhật, controller tra cứu trạng thái hiện tại của bản ghi
từ cơ sở dữ liệu, sau đó kiểm tra trạng thái mới có thuộc tập hợp cho phép
hay không. Nếu không hợp lệ, controller trả về HTTP 400 với thông báo lỗi cụ
thể trước khi thực hiện bất kỳ thao tác ghi nào.

\textbf{Cơ chế 2: Kiểm tra quyền theo trạng thái đích.} Không phải mọi tác
nhân đều có quyền thực hiện mọi bước chuyển đổi. Cụ thể, chỉ nhân viên hoặc
quản lý \break (\texttt{role = staff | manager}) được phép chuyển sang các trạng thái
từ \break \texttt{CONTACTING} đến \texttt{ITEM\_RECEIVED}; chỉ quản trị viên hoặc
quản lý 
\break (\texttt{role = admin | manager}) được phép chuyển sang \texttt{APPROVED}, \break
\texttt{REJECTED} hoặc \texttt{REFUNDED}. Middleware \texttt{verifyRole} kiểm
tra điều kiện này tại tầng route trước khi yêu cầu đến controller.

\subsubsection*{Thiết kế collection return\_request}

Collection \texttt{return\_request} lưu toàn bộ lịch sử chuyển đổi trạng thái
thông qua trường \texttt{status\_history} --- một mảng tài liệu nhúng, mỗi
phần tử ghi nhận trạng thái mới, email nhân viên thực hiện, thời điểm và ghi
chú. Thiết kế này cho phép hiển thị timeline đầy đủ trên giao diện quản lý mà
không cần truy vấn bổ sung sang collection khác. Trường \texttt{modified\_by}
lưu email nhân viên xử lý cuối cùng, hỗ trợ tra cứu nhanh trách nhiệm tại
bước hiện tại. Khi trạng thái chuyển sang \texttt{APPROVED} hoặc
\texttt{REJECTED}, hệ thống tự động gọi \texttt{Nodemailer} để gửi email thông
báo kết quả cho khách hàng, bao gồm lý do từ chối trong trường hợp yêu cầu
không được chấp nhận.

\subsection{Kết quả đạt được}
\label{subsection:5.1.3}

Giải pháp FSM cho quy trình hoàn trả hàng mang lại ba kết quả cụ thể. Thứ
nhất, toàn bộ 21 trường hợp kiểm thử liên quan đến hoàn hàng (bao gồm kiểm
thử vai trò và kiểm thử chuyển đổi trạng thái không hợp lệ) đều cho kết quả
đúng, xác nhận rằng hệ thống không cho phép bất kỳ chuyển đổi nào nằm ngoài
tập hợp đã định nghĩa. Thứ hai, giao diện quản lý hiển thị đúng danh sách yêu
cầu theo từng trạng thái với bộ lọc, cho phép nhân viên tập trung vào đúng
nhóm yêu cầu thuộc trách nhiệm của mình. Thứ ba, cơ chế ghi \texttt{status\_history}
cung cấp bằng chứng kiểm toán đầy đủ cho mỗi yêu cầu hoàn trả, đáp ứng nhu
cầu truy vết trách nhiệm trong vận hành thực tế.

%%%%%%%%%%%%%%%%%%%%%%%%%%%%%%%%%%%%%%%%%%
\section{Hệ thống phân quyền theo vai trò kết hợp nhật ký tác động hệ thống}
\label{section:5.2}

\subsection{Bài toán}
\label{subsection:5.2.1}

Khi một cửa hàng thời trang trực tuyến mở rộng đội ngũ nhân viên, việc quản lý
quyền truy cập hệ thống trở thành vấn đề nghiêm trọng. Nếu mọi nhân viên đều
có cùng quyền hạn, nguy cơ thao tác nhầm dữ liệu quan trọng (xóa sản phẩm,
thay đổi giá, xử lý đơn hàng không đúng thẩm quyền) là rất cao. Ngược lại,
nếu chỉ có quản trị viên mới được thao tác tất cả, hệ thống tạo ra nút thắt
cổ chai và không phân công được công việc theo nhóm.

Vấn đề bổ sung là ngay cả khi đã phân quyền, quản trị viên vẫn cần có khả
năng truy vết: ai đã thực hiện hành động gì, trên đối tượng nào, vào thời
điểm nào. Không có cơ chế này, khi xảy ra sự cố (đơn hàng bị thay đổi trạng
thái sai, sản phẩm bị xóa nhầm), quản trị viên không có căn cứ để điều tra.
Các nền tảng SaaS hiện tại như Haravan hay Sapo cung cấp phân quyền theo vai
trò ở mức cơ bản nhưng thiếu hoàn toàn tính năng nhật ký tác động hệ thống,
như đã xác định ở mục \ref{section:2.1}.

Bài toán đặt ra là xây dựng một hệ thống phân quyền đáp ứng đồng thời ba tiêu
chí: (i) kiểm soát truy cập tại tầng API, không chỉ ở tầng giao diện, để
ngăn chặn cả các cuộc tấn công trực tiếp vào API; (ii) cho phép cấu hình
quyền linh hoạt theo vai trò mà không cần thay đổi mã nguồn khi thêm nhân
viên mới hoặc thay đổi phân công; và (iii) ghi nhận toàn bộ hành động có tác
động đến dữ liệu vào nhật ký tác động hệ thống có thể tra cứu theo nhiều chiều.

\subsection{Giải pháp}
\label{subsection:5.2.2}

\subsubsection*{Mô hình RBAC hai tầng}

Giải pháp hiện thực mô hình \textbf{RBAC} (Role-Based Access Control) theo hai
tầng kiểm soát phân biệt.

\textbf{Tầng 1 --- Kiểm soát truy cập API (Middleware).} Mọi route trong nhóm
\texttt{/api/admin} đều đi qua chuỗi middleware hai bước:
\[
  \texttt{verifyToken} \;\longrightarrow\; \texttt{verifyRole}
\]
\texttt{verifyToken} giải mã JWT từ cookie và xác thực chữ ký; nếu token hết
hạn hoặc không hợp lệ, trả về HTTP 401 ngay lập tức. \texttt{verifyRole} tra
cứu vai trò của người dùng từ collection \texttt{role} trong MongoDB và so
sánh với danh sách \texttt{allowedActions} của route tương ứng; nếu không đủ
quyền, trả về HTTP 403. Kiểm soát này diễn ra hoàn toàn ở tầng server, độc
lập với giao diện, nên không thể bị vượt qua bằng cách tắt JavaScript hay
thao tác DOM trực tiếp trên trình duyệt.

\textbf{Tầng 2 --- Kiểm soát hiển thị giao diện (PermissionContext).} Phía
client, \texttt{PermissionContext} (mục \ref{subsection:4.1.2}) giữ danh sách
quyền hành động và danh sách menu được phép của người dùng đang đăng nhập.
Các component \texttt{SlideBar} và \texttt{HomePageAdmin} đọc context này để
chỉ render các mục menu mà người dùng có quyền, tránh hiển thị nhầm tính năng
không thuộc phạm vi vai trò. Tầng này là lớp trải nghiệm người dùng, bổ sung
cho tầng server nhưng không thay thế được tầng đó.

\subsubsection*{Cấu hình quyền động qua collection role}

Quyền của mỗi vai trò được lưu trong collection \texttt{role} với hai trường
chính: \texttt{permissions} (mảng mã hành động, ví dụ \texttt{EDIT\_PRODUCT}, \break
\texttt{UPDATE\_ORDER\_STATUS}) và \texttt{allowedMenus} (mảng tên route menu
được phép). Khi quản trị viên cần thay đổi quyền của một vai trò, thao tác
được thực hiện hoàn toàn qua giao diện quản lý (hình \ref{fig:screen-admin-role}),
không cần can thiệp vào mã nguồn. \break \texttt{verifyRole} luôn tra cứu quyền từ
cơ sở dữ liệu tại thời điểm xử lý request, đảm bảo thay đổi quyền có hiệu
lực ngay lập tức với request tiếp theo mà không cần khởi động lại server.

\subsubsection*{Nhật ký kiểm toán với tham chiếu đa hình}

Nhật ký kiểm toán được lưu trong collection \texttt{audit\_log} với thiết kế
\textbf{tham chiếu đa hình} (polymorphic reference): mỗi bản ghi chứa các
trường \break \texttt{target\_type} (loại đối tượng bị tác động, ví dụ
\texttt{"Product"}, \texttt{"Order"}, \break \texttt{"ReturnRequest"}) và
\texttt{target\_id} (ObjectId của đối tượng). Thiết kế này cho phép một
collection duy nhất ghi nhận hành động trên bất kỳ loại đối tượng nào trong
hệ thống, không cần tạo collection nhật ký riêng cho mỗi loại đối tượng.

Mỗi bản ghi nhật ký bao gồm: \texttt{actor\_email} (người thực hiện), \break
\texttt{action\_code} (mã hành động chuẩn hóa), \texttt{target\_type} và
\texttt{target\_id} (đối tượng bị tác động), \texttt{description} (mô tả
ngôn ngữ tự nhiên), và \texttt{created\_at} (thời điểm). Collection được
lập chỉ mục tổng hợp trên \texttt{(actor\_email, created\_at)},
\texttt{(target\_type, target\_id, created\_at)} và \break \texttt{action\_code}
để hỗ trợ các truy vấn nhật ký theo nhiều chiều: theo nhân viên, theo đối
tượng cụ thể, hoặc theo loại hành động.

Hàm tiện ích \texttt{auditLog(actor, action, targetType, \break targetId, description)}
được gọi từ controller sau mỗi thao tác có tác động đến dữ liệu, đảm bảo
không bỏ sót bất kỳ hành động nào cần ghi nhận.

\subsection{Kết quả đạt được}
\label{subsection:5.2.3}

Hệ thống phân quyền kết hợp nhật ký kiểm toán mang lại kết quả trên cả hai
phương diện bảo mật và vận hành. Về bảo mật, trường hợp kiểm thử TC-19
(mục \ref{section:4.4}) xác nhận nhân viên có vai trò \texttt{staff} không
thể phê duyệt yêu cầu hoàn hàng dù gọi thẳng API, hệ thống trả về HTTP 403
đúng như thiết kế. Về vận hành, quản trị viên có thể thay đổi toàn bộ cấu
hình quyền của một vai trò qua giao diện web trong vòng dưới một phút, không
cần can thiệp kỹ thuật. Nhật ký kiểm toán ghi nhận đầy đủ các hành động trên
sản phẩm, đơn hàng và yêu cầu hoàn trả, cho phép tra cứu lịch sử thao tác
theo nhân viên hoặc theo đối tượng bất kỳ thông qua trang lịch sử kiểm toán
trong giao diện quản trị.

%%%%%%%%%%%%%%%%%%%%%%%%%%%%%%%%%%%%%%%%%%
\section{Cấu hình phí vận chuyển linh hoạt}
\label{section:5.3}

\subsection{Bài toán}
\label{subsection:5.3.1}

Phí vận chuyển là một trong những yếu tố ảnh hưởng trực tiếp đến quyết định
mua hàng và biên lợi nhuận của cửa hàng. Đối với doanh nghiệp bán lẻ thời
trang có hợp tác với nhiều đơn vị vận chuyển khác nhau ở các khu vực địa lý
khác nhau, phí vận chuyển thực tế đến từng địa bàn là không đồng nhất: giao
hàng nội thành và ngoại thành có chi phí khác nhau; giao đến các huyện vùng
sâu có thể cao hơn đáng kể so với trung tâm thành phố; một số xã/phường cụ
thể có thể được miễn phí vận chuyển theo chính sách khuyến mãi tạm thời.

Các nền tảng thương mại điện tử hiện tại, bao gồm cả Haravan và \break WooCommerce,
thường chỉ hỗ trợ cấu hình phí vận chuyển ở cấp tỉnh/thành phố hoặc cấp
quận/huyện. Không có nền tảng nào trong nhóm khảo sát (bảng
\ref{tab:comparison}) hỗ trợ cấu hình đến cấp xã/phường. Điều này buộc doanh
nghiệp phải áp dụng mức phí trung bình hoặc cao nhất cho toàn quận/huyện, dẫn
đến mất cạnh tranh ở các địa bàn có chi phí thực tế thấp hơn, hoặc thua lỗ ở
những địa bàn chi phí cao hơn mức trung bình được cấu hình.

Bài toán cụ thể là thiết kế một cơ chế cho phép quản trị viên cấu hình mức
phí vận chuyển độc lập cho từng cấp địa bàn (tỉnh/thành $\to$ quận/huyện $\to$
xã/phường), và hệ thống tự động tính phí chính xác ngay tại trang thanh toán
dựa trên địa chỉ khách hàng nhập vào, không cần can thiệp thủ công.

\subsection{Giải pháp}
\label{subsection:5.3.2}

\subsubsection*{Thiết kế lưu trữ phân cấp địa lý trong MongoDB}

Thay vì sử dụng bảng phẳng với một hàng cho mỗi địa bàn (dẫn đến hàng nghìn
hàng khi cấu hình đến cấp xã/phường), giải pháp tận dụng khả năng lưu tài
liệu lồng nhau của MongoDB để biểu diễn cấu trúc phân cấp địa lý một cách
tự nhiên.

Collection \texttt{shipping\_config} lưu mỗi tài liệu tương ứng với một
tỉnh/thành phố. Bên trong mỗi tài liệu, trường \texttt{districts} là mảng
các tài liệu lồng nhau --- mỗi phần tử đại diện cho một quận/huyện và chứa
trường \texttt{fee} (phí mặc định cho quận đó) cùng trường \texttt{wards}
là mảng lồng tiếp theo, mỗi phần tử đại diện cho một xã/phường với
\texttt{fee} riêng. Cấu trúc này cho phép truy vấn theo chuỗi:
\begin{enumerate}
  \item Tìm tài liệu tỉnh/thành theo tên $\to$ trả về toàn bộ cấu hình quận/huyện.
  \item Trong tài liệu đó, lọc phần tử \texttt{districts} có tên khớp quận/huyện
        $\to$ trả về phí quận và danh sách xã/phường.
  \item Trong mảng \texttt{wards} của quận đó, tìm xã/phường theo tên $\to$
        trả về phí xã/phường nếu được cấu hình riêng.
\end{enumerate}

\textbf{Quy tắc ưu tiên:} Nếu xã/phường được cấu hình phí riêng, ưu tiên sử
dụng phí đó. Nếu không, sử dụng phí của quận/huyện. Nếu quận/huyện cũng chưa
được cấu hình, sử dụng phí mặc định toàn hệ thống (\texttt{DEFAULT\_SHIPPING\_FEE}).
Quy tắc ưu tiên này cho phép quản trị viên chỉ cần cấu hình những địa bàn có
phí khác biệt, không cần nhập đầy đủ tất cả xã/phường.

\subsubsection*{Tính phí động tại tầng API}

Tại tầng server, hàm \texttt{calculateShippingFee(province, \break district, ward)}
trong \texttt{ControllerShipping} thực hiện truy vấn MongoDB theo logic ưu
tiên nêu trên và trả về mức phí áp dụng. Hàm này được gọi ở hai điểm:
(i)~khi client gửi yêu cầu tính phí real-time tại trang thanh toán (endpoint
\texttt{GET /api/shipping-fee}) và (ii)~trong luồng tạo đơn hàng (endpoint
\texttt{POST /api/payment}) để ghi nhận phí vận chuyển vào bản ghi đơn hàng,
đảm bảo phí được tính tại thời điểm đặt hàng phản ánh đúng cấu hình hiện tại.

\subsubsection*{Component CascadingAddressSelect phía client}

Phía giao diện, component \texttt{CascadingAddressSelect} hiện thực bộ chọn
địa chỉ phân cấp dạng tầng (cascading). Dữ liệu danh sách tỉnh/thành phố,
quận/huyện và xã/phường của Việt Nam được tải từ một API địa lý công khai
(\texttt{provinces.open-api.vn}) và lưu vào \texttt{sessionStorage} để tránh
tải lại không cần thiết khi người dùng điều hướng giữa các trang. Khi người
dùng chọn xong địa chỉ cấp xã/phường, component tự động gọi endpoint tính
phí vận chuyển và cập nhật tổng tiền hiển thị trong tóm tắt giỏ hàng mà không
cần người dùng nhấn thêm bất kỳ nút nào.

\subsection{Kết quả đạt được}
\label{subsection:5.3.3}

Giải pháp cấu hình phí vận chuyển phân cấp mang lại hai kết quả chính. Về
tính năng, hệ thống cho phép quản trị viên thiết lập mức phí độc lập cho từng
xã/phường tùy ý, tức là độ phân giải cấu hình cao hơn bất kỳ nền tảng nào
trong nhóm khảo sát. Về trải nghiệm người dùng, phí vận chuyển được cập nhật
tức thời khi khách hàng hoàn thành chọn địa chỉ, cho phép khách hàng thấy
ngay tổng chi phí trước khi xác nhận đặt hàng, giảm thiểu nguy cơ hủy đơn
do bất ngờ về phí vận chuyển ở bước cuối. Kết quả kiểm thử triển khai (mục
\ref{section:4.5}) xác nhận thời gian phản hồi của endpoint tính phí vận
chuyển nằm trong ngưỡng dưới 200ms, đảm bảo trải nghiệm cập nhật real-time
mượt mà cho người dùng cuối.

%%%%%%%%%%%%%%%%%%%%%%%%%%%%%%%%%%%%%%%%%%
\section{Tích hợp thanh toán VNPay theo luồng redirect an toàn}
\label{section:5.4}

\subsection{Bài toán}
\label{subsection:5.4.1}

Tích hợp cổng thanh toán trực tuyến là yêu cầu bắt buộc của hệ thống thương
mại điện tử, nhưng cũng là điểm dễ mắc lỗi bảo mật nhất. Hai rủi ro chính
trong tích hợp VNPay là: (i) giả mạo callback --- kẻ tấn công gọi trực tiếp
endpoint callback với tham số giả mạo để đánh dấu đơn hàng là đã thanh toán
mà không thực sự chuyển tiền; và (ii) sửa đổi tham số --- kẻ tấn công thay
đổi số tiền hoặc mã đơn hàng trong URL được tạo ra trước khi chuyển hướng
đến cổng VNPay. Nếu không xử lý đúng, cả hai lỗi này đều dẫn đến thiệt hại
tài chính trực tiếp cho cửa hàng.

\subsection{Giải pháp}
\label{subsection:5.4.2}

Giải pháp hiện thực luồng VNPay hoàn toàn theo tài liệu kỹ thuật chính thức
của VNPay với hai cơ chế bảo vệ bổ sung.

\textbf{Ký số URL thanh toán bằng HMAC-SHA512.} Khi tạo URL chuyển hướng đến
VNPay, server tính chữ ký HMAC-SHA512 trên chuỗi tham số đã được sắp xếp
theo thứ tự cố định (bao gồm số tiền, mã đơn hàng, thời điểm và mã merchant),
sử dụng khóa bí mật lưu trong biến môi trường \texttt{VNP\_HASH\_SECRET}.
Chữ ký được gắn vào URL dưới dạng tham số \texttt{vnp\_SecureHash}. Cổng
VNPay sẽ từ chối mọi yêu cầu có chữ ký không khớp, ngăn chặn sửa đổi tham
số.

\textbf{Xác thực chữ ký callback.} Khi VNPay gọi callback URL \break \texttt{GET
/vnpay-return}, server không tin tưởng bất kỳ tham số nào ngay lập tức. Thay
vào đó, server tách \texttt{vnp\_SecureHash} ra khỏi danh sách tham số, sắp
xếp lại các tham số còn lại theo đúng thứ tự, tính lại chữ ký HMAC-SHA512 và
so sánh với giá trị nhận được. Chỉ khi hai chữ ký khớp, server mới cập nhật \break
\texttt{statusPayment = true} trong \texttt{ModelOrder}. Cơ chế này loại bỏ
hoàn toàn khả năng giả mạo callback vì kẻ tấn công không có khóa bí mật để
tạo chữ ký hợp lệ.

\subsection{Kết quả đạt được}
\label{subsection:5.4.3}

Luồng thanh toán VNPay sandbox được kiểm thử với tất cả các mã phản hồi do
VNPay cung cấp trong môi trường thử nghiệm (mục \ref{section:4.5}). Các
trường hợp kiểm thử TC-12, TC-13 và TC-14 (mục \ref{section:4.4}) xác nhận
hệ thống xử lý đúng cả ba kịch bản: chuyển hướng thành công, callback thành
công và callback thất bại. Đặc biệt, cơ chế xác thực chữ ký ngăn chặn việc
đơn hàng bị đánh dấu thanh toán trong trường hợp callback giả mạo, đảm bảo
tính toàn vẹn của dữ liệu tài chính.

\end{document}
